\lesson{4}{Add a damper}{A force against velocity drains the energy the spring cannot — and one number, ζ, says how much.}{1}{damper}{Part I · PID}
\label{les:damper}

\lead{Lesson 2 left us with a drone that bounced forever. The cure is the third letter of PID. The derivative term looks at how fast the error is changing and pushes against it — the shock absorber to the proportional spring. In this lesson we solve the damped equation completely, derive the formulas that engineers use to tune it, and check every one of them against the simulator.}

\section{The derivative term is a damper}

With a constant setpoint, the error changes only because the drone moves: $\dot e = \dot r - \dot y = -\dot y$. The derivative term is therefore
\begin{equation}
  D = K_d\,\dot e = -K_d\,\dot y ,
\end{equation}
a force proportional to the vertical velocity and opposing it — exactly what a shock absorber does. With exact feedforward, P and D together give the closed loop (with $x = y - r$, as in \lessonref{les:spring})
\begin{equation}\label{eq:damper-msd}
  m\,\ddot x + K_d\,\dot x + K_p\,x = 0 ,
\end{equation}
the equation of a \term{mass–spring–damper}.

\subsection{The damper removes energy}
Repeat the energy calculation of \lessonref{les:spring}: multiply \cref{eq:damper-msd} by $\dot x$ and use the chain rule,
\begin{equation}
  \frac{\dd}{\dd t}\Big(\tfrac12 m\dot x^2 + \tfrac12 K_p x^2\Big) = \dot x\,(m\ddot x + K_p x) = -K_d\,\dot x^2 \;\le\; 0 .
\end{equation}
Whenever the drone moves, the D term removes energy at the rate $K_d\dot x^2$ — the power dissipated by a damper. The energy can only decrease, and it stops decreasing only when the drone stops moving; since a spring that is stretched does not let the drone stay still, the only final state is $x = 0$, $\dot x = 0$. The oscillation must die out. The remaining questions are \emph{how fast} and \emph{in what manner}.

\section{Solving the damped equation}

We follow the recipe of \lessonref{les:spring} step by step.

\paragraph{The characteristic equation.} Substituting $x = e^{st}$ into \cref{eq:damper-msd} gives $m s^2 + K_d s + K_p = 0$. Divide by $m$ and introduce two new parameters:
\begin{equation}\label{eq:damper-std}
  s^2 + 2\zeta\omega_n\,s + \omega_n^2 = 0,\qquad
  \omega_n = \sqrt{\frac{K_p}{m}},\qquad
  \zeta = \frac{K_d}{2\sqrt{K_p\,m}} .
\end{equation}
(Check: $2\zeta\omega_n = 2\cdot\frac{K_d}{2\sqrt{K_pm}}\cdot\sqrt{\frac{K_p}{m}} = \frac{K_d}{m}$ and $\omega_n^2 = \frac{K_p}{m}$, as required.) This \emph{standard form} separates two questions. The \term{natural frequency} $\omega_n$ sets the time scale: how fast the loop wants to move. The \term{damping ratio} $\zeta$ (zeta) is dimensionless and sets the shape: how the motion approaches the target.

\paragraph{The roots.} The quadratic formula gives
\begin{equation}\label{eq:damper-poles}
  s_{1,2} = -\zeta\omega_n \pm \omega_n\sqrt{\zeta^2 - 1} .
\end{equation}
The sign of $\zeta^2 - 1$ decides everything: three cases follow.

\begin{definition}{The three kinds of damping}
\begin{description}[leftmargin=0pt, itemsep=3pt, font=\sffamily\small\bfseries]
  \item[Underdamped, $0<\zeta<1$.] The square root is imaginary: complex poles $-\zeta\omega_n \pm j\omega_d$ with the \emph{damped frequency} $\omega_d = \omega_n\sqrt{1-\zeta^2}$. The motion is a decaying oscillation, $x(t) = e^{-\zeta\omega_n t}(A\cos\omega_d t + B\sin\omega_d t)$. The response overshoots.
  \item[Critically damped, $\zeta = 1$.] The root is zero: a double real pole at $-\omega_n$, and $x(t) = (A + Bt)\,e^{-\omega_n t}$. The fastest approach that never overshoots.
  \item[Overdamped, $\zeta > 1$.] Two distinct real poles, $x(t) = A e^{s_1 t} + B e^{s_2 t}$. The one closer to the origin makes the response slow: the damper is so strong that it resists the motion towards the target too.
\end{description}
\end{definition}

\widefig{zeta_family}{The damping ratio in two pictures. Left: the poles of \cref{eq:damper-std}, scaled by $\omega_n$. For $\zeta<1$ they lie on the unit circle (their distance from the origin is $\omega_n$) and move towards the real axis as $\zeta$ grows; at $\zeta=1$ they meet; beyond it they split along the axis and one of them crawls towards the origin. Right: the corresponding step responses. The grey band is the $\pm2\,\%$ tolerance used for the settling time.}{fig:zeta-family}

\Cref{fig:zeta-family} is worth a long look. The pole plane is a map of behaviour: the height of a pole above the real axis is the frequency of oscillation, its distance to the left of the imaginary axis is the speed of decay, and — since $|s| = \omega_n$ for complex poles — the angle $\theta$ from the negative real axis tells the damping, $\cos\theta = \zeta$.

\begin{example}[The step response of the underdamped loop]
\label{ex:damper-under}
The setpoint jumps by \SI{1}{m} and $K_d = 3$ ($K_p = 10$, $m = 1$). Find the motion.

\textbf{Step 1 — the parameters.} $\omega_n = \sqrt{10} = \SI{3.162}{rad/s}$, $\zeta = 3/(2\sqrt{10}) = 0.474$, $\zeta\omega_n = K_d/(2m) = \SI{1.5}{s^{-1}}$, $\omega_d = 3.162\sqrt{1 - 0.225} = \SI{2.784}{rad/s}$.

\textbf{Step 2 — the general solution.} $x(t) = e^{-1.5t}(A\cos 2.784t + B\sin 2.784t)$.

\textbf{Step 3 — initial conditions.} Just after the step, $x(0) = -1$ (one metre below the new setpoint) and $\dot x(0) = 0$. The first gives $A = -1$. Differentiating, $\dot x(0) = -1.5A + 2.784B = 0$, so $B = 1.5A/2.784 = -0.539$.

\textbf{Step 4 — the result.} $x(t) = -e^{-1.5t}\big(\cos 2.784t + 0.539\sin 2.784t\big)$.

\textbf{Step 5 — interpret.} The oscillation has the period $2\pi/\omega_d = \SI{2.26}{s}$ — slightly longer than the undamped \SI{1.99}{s} — and its envelope shrinks like $e^{-1.5t}$: to half in $\ln 2/1.5 = \SI{0.46}{s}$.
\end{example}

\section{Reading a step response}

To compare controllers we need numbers, not adjectives. The simulator measures four of them after every setpoint step, and shows them in the \ui{step metrics} table.

\begin{definition}{Step-response metrics}
For a step of the setpoint from $y_0$ to $y_1$:
\begin{description}[leftmargin=0pt, itemsep=2pt, font=\sffamily\small\bfseries]
  \item[Rise time $t_r$] from 10\,\% to 90\,\% of the step.
  \item[Overshoot $M_p$] the peak beyond $y_1$, in \% of the step size.
  \item[Settling time $t_s$] the time after which the output stays within $\pm2\,\%$ of the step around $y_1$.
  \item[Steady-state error $e_{ss}$] what is left once everything has settled.
\end{description}
\end{definition}

For the second-order loop these metrics follow from the solution. We derive the two most used ones.

\begin{example}[Deriving the overshoot formula]
\label{ex:damper-mp}
Show that the overshoot of an underdamped second-order loop after a unit step is $M_p = e^{-\pi\zeta/\sqrt{1-\zeta^2}}$.

\textbf{Step 1 — the unit step response.} Repeating Example~\ref{ex:damper-under} with symbols ($x(0) = -1$, $\dot x(0) = 0$) gives, for $y = 1 + x$,
\begin{equation*}
  y(t) = 1 - e^{-\zeta\omega_n t}\Big(\cos\omega_d t + \tfrac{\zeta}{\sqrt{1-\zeta^2}}\sin\omega_d t\Big).
\end{equation*}

\textbf{Step 2 — the velocity.} Differentiating and simplifying (the cosine terms cancel) leaves $\dot y(t) = \frac{\omega_n}{\sqrt{1-\zeta^2}}\,e^{-\zeta\omega_n t}\sin\omega_d t$.

\textbf{Step 3 — the peak.} The first maximum is where the velocity first returns to zero after $t = 0$: $\sin\omega_d t = 0$, i.e.\ $t_p = \pi/\omega_d$.

\textbf{Step 4 — evaluate.} At $t_p$: $\cos\pi = -1$, $\sin\pi = 0$, so $y(t_p) = 1 + e^{-\zeta\omega_n\pi/\omega_d} = 1 + e^{-\pi\zeta/\sqrt{1-\zeta^2}}$.

\textbf{Step 5 — the result.} The overshoot is the part above 1:
\begin{equation}\label{eq:damper-mp}
  M_p = e^{-\pi\zeta/\sqrt{1-\zeta^2}} .
\end{equation}
It depends on $\zeta$ only. For $\zeta = 0.474$: $M_p = e^{-\pi\cdot0.474/0.881} = e^{-1.69} = 0.184$, i.e.\ 18.4\,\%, reached at $t_p = \pi/2.784 = \SI{1.13}{s}$.
\end{example}

\paragraph{The settling time.} The oscillation in Step 1 is multiplied by the envelope $e^{-\zeta\omega_n t}$ (times a factor close to one). The response is inside the 2\,\% band once the envelope has fallen below 0.02, i.e.\ when $\zeta\omega_n t = \ln 50 \approx 3.9$. Rounding up gives the rule
\begin{equation}\label{eq:damper-ts}
  t_s \approx \frac{4}{\zeta\,\omega_n}.
\end{equation}
For $K_d = 3$: $t_s \approx 4/1.5 = \SI{2.7}{s}$.

\begin{example}[The critically damped loop]
\label{ex:damper-crit}
Find the $K_d$ that gives $\zeta = 1$, and the settling time of the resulting step response ($K_p = 10$, $m = 1$).

\textbf{Step 1 — the gain.} $\zeta = 1$ means $K_d = 2\sqrt{K_p m} = 2\sqrt{10} = 6.32$.

\textbf{Step 2 — the response.} With the double pole at $-\omega_n$, the general solution $x = (A + Bt)e^{-\omega_n t}$ with $x(0) = -1$, $\dot x(0) = 0$ gives $A = -1$, $B = -\omega_n$, so
$y(t) = 1 - (1 + \omega_n t)\,e^{-\omega_n t}$.

\textbf{Step 3 — settling.} We need $(1 + \omega_n t)e^{-\omega_n t} = 0.02$. With $z = \omega_n t$, solve $(1+z)e^{-z} = 0.02$ by iteration, $z \leftarrow \ln\big(50(1+z)\big)$: starting from $z = 4$ gives $5.52$, then $5.77$, $5.81$, $5.83$ — converged: $z = 5.83$.

\textbf{Step 4 — numbers.} $t_s = 5.83/3.162 = \SI{1.84}{s}$. The rise time (solving $y = 0.1$ and $y = 0.9$ the same way: $z = 0.53$ and $3.89$) is $(3.89 - 0.53)/3.162 = \SI{1.06}{s}$.
\end{example}

\begin{example}[The overdamped loop]
\label{ex:damper-over}
With $K_d = 15$ ($\zeta = 2.37$), where are the poles, and why is the response slow?

\textbf{Step 1 — the poles.} From \cref{eq:damper-poles}: $s = -2.37\cdot3.162 \pm 3.162\sqrt{2.37^2 - 1} = -7.50 \pm 6.80$, i.e.\ $s_1 = -0.70$ and $s_2 = -14.3$.

\textbf{Step 2 — who dominates.} Both modes decay, but $e^{-14.3t}$ is gone within \SI{0.3}{s}, while $e^{-0.70t}$ has a time constant of \SI{1.43}{s}. After the first fraction of a second, only the slow pole matters: it is the \term{dominant pole}.

\textbf{Step 3 — settling.} The response enters the 2\,\% band when $e^{-0.70t} \approx 0.02$: $t \approx \ln 50/0.70 = \SI{5.6}{s}$ — three times slower than the critically damped loop, with a damper more than twice as strong.
\end{example}

\section{The experiment}

The lesson flies a square wave: the setpoint jumps between \SI{1.5}{m} and \SI{2.5}{m} every four seconds. \Cref{fig:damper-sweep} overlays one upward step for five values of $K_d$ (with $K_p = 10$, $m=\SI{1}{kg}$, so $\omega_n = \SI{3.16}{rad/s}$ and $\zeta = K_d/6.32$).

\simfig{damper_sweep}{Five dampers, one step. The low-$\zeta$ runs start slightly off because the previous step had not finished settling.}{fig:damper-sweep}

\begin{table}[h]
\centering\small\sffamily
\begin{tabular}{rrrrrrr}
\toprule
 & & \multicolumn{2}{c}{overshoot} & \multicolumn{2}{c}{settling time} & rise time\\
\cmidrule(lr){3-4}\cmidrule(lr){5-6}
$K_d$ & $\zeta$ & theory & measured & theory & measured & th.\,/\,meas.\\
\midrule
3 & 0.47 & 18.4\,\% & 15.8\,\% & \SI{2.7}{s} & \SI{2.5}{s} & — / \SI{0.50}{s}\\
6.3 & 1.00 & 0 & 0 & \SI{1.84}{s} & \SI{1.97}{s} & \SI{1.06}{s} / \SI{1.09}{s}\\
15 & 2.37 & 0 & 0 & \SI{5.6}{s} & $>\SI{4}{s}$ & — / \SI{3.2}{s}\\
\bottomrule
\end{tabular}
\caption{Theory (Examples~\ref{ex:damper-mp}–\ref{ex:damper-over}) against the simulator. The small differences come from what \cref{eq:damper-msd} leaves out: the \SI{30}{ms} motor lag, sampling at \SI{250}{Hz}, the low-pass filter on D and the air drag.}
\label{tab:damper}
\end{table}

\Cref{tab:damper} puts the formulas next to the measurements. The agreement is good enough to design with — which is the point of a model. The lesson has a goal: overshoot below 2\,\% and settling below \SI{2.5}{s}. The critically damped $K_d \approx 6.3$ meets it.

\begin{keyidea}[Critical damping]
For a given stiffness the fastest response without overshoot is the critically damped one, $\zeta = 1$, i.e.\ $K_d = 2\sqrt{K_p m}$. For the default $K_p = 10$ that is $K_d \approx 6.3$. Many practical tunes aim slightly below, $\zeta \approx 0.7$–$0.9$: a few percent of overshoot buys a noticeably faster rise.
\end{keyidea}

\begin{example}[Designing for a given overshoot]
\label{ex:damper-design}
Which $K_d$ gives 5\,\% overshoot with $K_p = 10$, $m = 1$? How long does it take to settle?

\textbf{Step 1 — invert \cref{eq:damper-mp}.} $\ln M_p = -\pi\zeta/\sqrt{1-\zeta^2}$. With $L = -\ln 0.05 = 3.00$: $L^2(1 - \zeta^2) = \pi^2\zeta^2$, so $\zeta = L/\sqrt{\pi^2 + L^2} = 3.00/\sqrt{9.87 + 9.00} = 0.69$.

\textbf{Step 2 — the gain.} $K_d = 2\zeta\sqrt{K_pm} = 2\cdot0.69\cdot3.162 = 4.36$.

\textbf{Step 3 — settling.} $t_s \approx 4/(\zeta\omega_n) = 4/(0.69\cdot3.162) = \SI{1.83}{s}$ — about the same as critical damping, but with a faster rise.
\end{example}

The default tune of Anemostatos, $K_d = 7$, is $\zeta = 1.11$: slightly overdamped, deliberately, because the integral term — next lesson — adds some overshoot of its own.

\screen[0 0 495 998]{damper}{Lesson 4 in the app, with $K_d = 0.8$: a ringing square-wave response. The info card (right) says \enquote{underdamped — overshoots} and computes $\zeta$ from the gains; the tracking chart marks the peak and the settling point.}{fig:damper-screen}

\begin{inapp}[Step metrics and the goal]
\begin{itemize}[leftmargin=1.2em]
  \item After every setpoint step \src{src/engine/metrics.ts} (\texttt{analyzeLastStep}) measures rise time, overshoot, settling time and steady-state error from the telemetry, with the definitions above; \src{src/ui/charts/StepMetrics.tsx} shows them.
  \item The lesson has a \emph{goal}: overshoot below 2\,\% and settling time below \SI{2.5}{s}. The lesson panel turns green when the last step meets it.
  \item The square wave is \ui{Setpoint\then Automatic motion\then square}, with amplitude and period below it (\param{setpoint.profile}).
\end{itemize}
\end{inapp}

\begin{trythis}
\begin{enumerate}
  \item Raise $K_d$ step by step from 0.8: 2, 4, 6.3. Watch the overshoot and the settling time in the step metrics, and $\zeta$ in the info card.
  \item Go past critical: 15, then 30. The response gets lazy — and the settling time grows again.
  \item Set $K_d = 4.36$ and check the overshoot against Example~\ref{ex:damper-design}.
  \item Set $K_p = 40$ and find the $K_d$ that is critical for it. How much faster is the response?
\end{enumerate}
\predict{with $K_p = 40$ and critical damping, what settling time do you expect? (Use Example~\ref{ex:damper-crit}.)}
\end{trythis}

\begin{curiosity}{A 660-tonne pendulum}
\photo{taipei}{The tuned mass damper of Taipei 101: a steel sphere of 660 tonnes hanging between the 87th and the 92nd floor.}
Taipei 101 was the tallest building in the world when it opened in 2004, and it stands in a place that is hit by typhoons and earthquakes. A tall building is a spring with very little damping: steel and concrete dissipate almost no energy, so wind gusts can set it swaying in a way that makes the people at the top seasick.

The engineers hung a giant pendulum — a sphere of welded steel plates weighing 660 tonnes — near the top of the tower, tuned so that its natural frequency matches the building's. When the tower sways, the pendulum lags behind and pushes against the motion; eight large hydraulic dampers connected to it turn that relative motion into heat. In our language, the pendulum and its dampers add a large $K_d$ to a structure that had almost none, and they cut the swaying by up to 40\,\% in strong winds. During Typhoon Soudelor in 2015 the sphere swung by a metre.

Your car contains four much smaller dampers doing the same job. Car suspensions are usually tuned to $\zeta \approx 0.2$–$0.4$ — well below critical — because a completely non-oscillating ride feels harsh. Comfort, it turns out, is a choice of damping ratio.
\end{curiosity}

\begin{remember}
\begin{itemize}
  \item With a constant setpoint the D term is $-K_d\dot y$: a damper. It removes energy at the rate $K_d\dot y^2$.
  \item P and D make a mass–spring–damper with $\omega_n = \sqrt{K_p/m}$ and $\zeta = K_d/(2\sqrt{K_p m})$; its poles are $-\zeta\omega_n \pm \omega_n\sqrt{\zeta^2-1}$.
  \item $\zeta<1$ overshoots by $e^{-\pi\zeta/\sqrt{1-\zeta^2}}$, $\zeta=1$ is the fastest without overshoot, $\zeta>1$ crawls on its dominant pole.
  \item Settling time $\approx 4/(\zeta\omega_n)$ for $\zeta<1$, $5.83/\omega_n$ for $\zeta=1$.
\end{itemize}
\end{remember}

\begin{exercises}
\exercise[1] The drone carries a payload so that $m = \SI{1.5}{kg}$, with $K_p = 10$. Which $K_d$ is critical? If $K_d$ stays at 6.3, what are $\zeta$ and the overshoot?
\answer{$K_d = 2\sqrt{15} = 7.75$. With $K_d = 6.3$: $\zeta = 0.81$, $M_p \approx 1.3\,\%$.}
\exercise[1] Find $K_p$ and $K_d$ for $m = \SI{1}{kg}$ such that $\zeta = 0.7$ and $t_s \approx \SI{1}{s}$.
\answer{$\zeta\omega_n = 4$, $\omega_n = 5.71$, $K_p = 32.7$, $K_d = 8.0$.}
\exercise[1] For $K_d = 3$, compute the time of the first peak and the time of the first undershoot below the setpoint after it.
\answer{$t_p = \pi/\omega_d = \SI{1.13}{s}$; the next extremum at $2\pi/\omega_d = \SI{2.26}{s}$.}
\exercise[2] Show that for $\zeta > 1$ the slow pole satisfies $s \approx -\omega_n/(2\zeta)$ when $\zeta \gg 1$. Interpret: the drone then moves as if only the damper and the spring existed, with no mass.
\answer{$-\zeta\omega_n + \zeta\omega_n\sqrt{1-1/\zeta^2} \approx -\omega_n/(2\zeta) = -K_p/K_d$: the pole of $K_d\dot x + K_p x = 0$.}
\exercise[2] Using the solution of Example~\ref{ex:damper-mp}, show that the ratio of two successive peaks of the oscillation (measured from the setpoint) is $e^{-2\pi\zeta/\sqrt{1-\zeta^2}}$. How can you use this to measure $\zeta$ from a recorded response?
\answer{Successive maxima are $2\pi/\omega_d$ apart, so the envelope ratio is $e^{-\zeta\omega_n\cdot2\pi/\omega_d}$. Measuring the ratio $\rho$ of two peaks gives $\delta = -\ln\rho$ (logarithmic decrement) and $\zeta = \delta/\sqrt{4\pi^2 + \delta^2}$.}
\exercise[3] The pole angle: show that complex poles of \cref{eq:damper-std} make the angle $\theta$ with the negative real axis such that $\cos\theta = \zeta$. What region of the $s$-plane guarantees overshoot below 5\,\% and settling within \SI{2}{s}?
\answer{$|s| = \omega_n$, $\operatorname{Re}s = -\zeta\omega_n$, so $\cos\theta = \zeta$. Overshoot $\le 5\,\%$ needs $\zeta \ge 0.69$ ($\theta \le 46^\circ$); $t_s \le 2$ needs $\operatorname{Re}s \le -2$. The intersection of the wedge and the half-plane.}
\end{exercises}
